% This is samplepaper.tex, a sample chapter demonstrating the
% LLNCS macro package for Springer Computer Science proceedings;
% Version 2.21 of 2022/01/12

\documentclass[runningheads]{llncs}

\usepackage[T1]{fontenc}
\usepackage{graphicx}
\usepackage{todonotes}

\begin{document}

\title{\texttt{cc2jolie}: Certified Projection of Choreographies to Jolie}

% If the paper title is too long for the running head, you can set
% an abbreviated paper title here
%\titlerunning{Abbreviated paper title}

\author{
  Luís Cruz-Filipe \orcidID{0000-0002-7866-7484} \and
  Lovro Lugović \orcidID{0000-0001-9684-9567} \and
  Fabrizio Montesi \orcidID{0000-0003-4666-901X}
}

% First names are abbreviated in the running head.
% If there are more than two authors, 'et al.' is used.
\authorrunning{L. Cruz-Filipe et al.}

\institute{
  Department of Mathematics and Computer Science, University of Southern Denmark
  \email{\{lcf, lugovic, fmontesi\}@imada.sdu.dk}
}

\maketitle

\begin{abstract}
  % TODO: The abstract should briefly summarize the contents of the paper in 150--250 words.

  \emph{Choreographic programming} is an emerging programming paradigm whose goal is to help programmers write correct distributed programs, such as security protocol implementations and distributed systems \todo{Too general?}.
  Traditionally, the programmer is tasked with writing the implementation of each \emph{endpoint (role)} separately, having to take great care in matching all of the necessary send--receive communications.
  Such an approach has been demonstrated to be cumbersome and error-prone \cite{LLLG16}, which is why choreographic programming instead seeks to solve the problem from a global viewpoint.
  A \emph{choreography} is a specification of the overall behaviour of a distributed system, from which implementations of all of the involved endpoints can be generated automatically through a process known as \emph{endpoint projection (EPP)}.
  Because it's impossible to write unmatched communications, choreographies give us the benefit of deadlock freedom by design.
  Of course, this approach is only as correct as the choreographic compiler used, which is why there has been interest in formally verifying the correctness of such compilers.
  For the first time we build a choreographic compiler to Jolie, using a formally verified EPP procedure that has been extracted from the Coq formalization of CC. \todo{Coq formalization of?}

  \keywords{TODO: First keyword \and Second keyword \and Another keyword.}
\end{abstract}

\section{Introduction}

\begin{figure}[h]
  \centering
  \includegraphics{pipeline.drawio.pdf}
  \caption{The pipeline of our CC compiler.}
  \label{fig:pipeline}
\end{figure}

Test \cite{MGZ14}.

\subsubsection{Acknowledgements}

Please place your acknowledgments at the end of the paper, preceded by an unnumbered run-in heading (i.e. 3rd-level heading).

% ---- Bibliography ----
%
% BibTeX users should specify bibliography style 'splncs04'.
% References will then be sorted and formatted in the correct style.

\bibliographystyle{splncs04}
\bibliography{main}

\end{document}
